\section{Introduction}\label{sec:intro}

\subsection{The problem}\label{sec:problem}

The Clay problem on quantum Yang--Mills theory asks for a four-dimensional quantum gauge theory with compact simple gauge group that satisfies axioms at least as strong as those of Osterwalder and Schrader \cite{os1973,os1975}, and that has a mass gap \cite{jaffewitten}. One constructive approach starts from Wilson's lattice action on a torus with lattice spacing $\varepsilon$ and takes $\varepsilon\to0$. The bare coupling $g_0(\varepsilon)$ must then go to zero at the rate that asymptotic freedom prescribes, so that the effective coupling at a fixed physical scale stays fixed. This paper treats the first limit only: lattice spacing to zero on a fixed torus. It does not treat the infinite-volume limit or the mass gap (Section~\ref{sec:notgiven}).

\emph{Identifiers.} Each object of the paper has an identifier of the form P-nnn. The prefix gives the kind: R for an input statement (most of them Ba{\l}aban-level), H for a hypothesis, condition or reading, L for a result proved here, K for a convention, Q for a closure item, X for an external result. At its first mention in a subsection, an object is named and its identifier follows in parentheses, for example ``the inequality (H2) (R-025)''. Later mentions use the name only. The registry \file{IDS.md} in the repository \cite{ymrepo} lists every identifier. The dependency table (Section~\ref{sec:deptable}), the table of inputs (Section~\ref{sec:rem-ledger}) and the appendices use the identifiers as their index.

\subsection{Ba{\l}aban's programme and the missing Theorem 2}\label{sec:balaban}

Between 1984 and 1989 Ba{\l}aban built a renormalization-group (RG) analysis of lattice gauge theories \cite{balaban95,balaban96,balaban98,balaban99a,balaban99b,balaban102a,balaban102b,balaban109,balaban116,balaban119,balaban122a,balaban122b}. His main output for $d=4$ is Theorem 1 of CMP 122 II: uniform bounds on the effective densities at every RG scale, under the hypothesis that the sequence of effective couplings stays small. His Theorem 2, which controls the running coupling, was to discharge this hypothesis. The two statements read as follows.

Theorem 2 \cite[p.~259]{balaban109}: ``Let $d = 4$, $G = SU(2)$, and let $\gamma$ be a sufficiently small positive constant, then for a sufficiently small positive $g$ there exists a bare coupling constant $g_0 = g_0(\varepsilon, g)$ such that the sequence of the effective coupling constants $g_k$ is contained in the interval $]0, \gamma]$, and $g_K = g$. Moreover, there exist constants $\beta, \beta'$, $0 < \beta \le \beta'$, such that $g_k$ satisfy the inequality $1/g^2 + \beta \log(L^k\varepsilon)^{-1} \le 1/g_k^2 \le 1/g^2 + \beta' \log(L^k\varepsilon)^{-1}$. (0.31)'' Ba{\l}aban adds that the proof ``will be given in a separate paper''. It was never published. In \cite[p.~355]{balaban122b} he writes that it ``has not been published yet, so we have Theorem 1 with the assumption''.

Theorem 1 \cite[p.~355]{balaban122b}: ``If the sequence of the effective coupling constants is contained in an interval $]0, \gamma]$ with a sufficiently small positive $\gamma$, then the effective densities $\rho_k$ have the form, and satisfy all the conditions and bounds, described in Sect.\ 2 [III].'' The bound (0.1) there reads $\chi_k \exp[-(1/g_k^2)A(U_k) - E_-|T_\eta|] \le \rho_k \le \exp E_+|T_\eta|$.

Without Theorem 2, the hypothesis of Theorem 1 is an assumption on the unknown coupling flow. Every continuum statement built on Ba{\l}aban's bounds inherits this assumption.

\subsection{The two targets and the main results}\label{sec:main}

The paper has two targets. The first target is a proof of enough of Theorem 2 to remove the coupling hypothesis from Theorem 1. The second target is a compactness theorem for the lattice Schwinger functions on a fixed torus, with the Osterwalder--Schrader properties proved for every limit point. Both results are conditional: each holds under named hypotheses, stated in the theorem.

\begin{mainthm}[weak Theorem 2]\label{thm:L002}
Let $G=SU(2)$, $d=4$, and $L$ odd, $L>11$. (For odd $L\ge3$ the statement holds if Ba{\l}aban's inductive assumption (H0) (H-202) and his construction of $\beta_j(g)$ at this $L$ are added as hypotheses.)
\begin{itemize}
\item[(i)] Assume the following.
\begin{itemize}
\item The uniform inequality (H2) (R-025) of Section~\ref{sec:twostep}.
\item The fixed-level inputs (PI) (R-026), (F) (R-027), ($M_j$) (R-028), V1 (R-029), and (G2) (R-030) of Section~\ref{sec:routeK}.
\item The history condition L5 (R-031) of Section~\ref{sec:L5}.
\item The bound $|\beta_{k+1}(h,g)-\beta_{k+1}(0)|\le C_\beta g$ on every past $h$, including non-monotone pasts (item 7 of Section~\ref{sec:remains}).
\end{itemize}
Then the conclusions (b), (c-pair), and K-018 (S) of Theorem 2w (Theorem~\ref{thm:L001}) hold.
\item[(ii)] Assume the conclusions of Theorem 2w, the condition L5, and the volume hypothesis (Vol) (K-032) of Section~\ref{sec:lattice}. Then Ba{\l}aban's Theorem 1 holds for the bare coupling $g_0(K,g)$, for every $K$ and every $g\in\,]0,g_1]$, without his coupling assumption.
\end{itemize}
\end{mainthm}
Part (i) follows from Theorem K (Theorem~\ref{thm:L009}) and Lemma DT (Lemma~\ref{lem:L005}). Part (ii) uses Ba{\l}aban's published proof of Theorem 1, the cross-scale sweep of Section~\ref{sec:sweep}, the list of later constants (Proposition~\ref{prop:L023}), and the finite-volume identification of Proposition~\ref{prop:L027} (Section~\ref{sec:conclusion}). A second route (Route B, \ThmPW{} in Section~\ref{sec:routeB}) also gives (i), under its own hypotheses (items B1--B7 of Section~\ref{sec:rem-inputs}). Both routes rest on (H2).

The hypotheses of (i) hold in the following sense.
\begin{itemize}
\item (H2) holds at $L=13$ (Theorem~\ref{thm:L046}) and at $L=15,17,21$ (Proposition~\ref{prop:H2other} in Section~\ref{sec:numr7}). At other $L$ it is a hypothesis.
\item (F) holds by Theorem~\ref{thm:L041}, ($M_j$) by Theorem~\ref{thm:L029}, and L5 by Theorem~\ref{thm:L030}, each relative to the inputs listed there.
\item (PI), (G2) and V1 hold at each fixed level, relative to the inputs listed in Section~\ref{sec:src-T1}.
\item The bound of item 7 holds under (H0) (R-006), or under the inputs of (G-gain) (H-111) (R-037).
\end{itemize}
So at $L=13$ the weak form of Theorem 2 rests on the inputs of the fixed-level results, of L5 and of item 7, on the Ba{\l}aban-level statements of Section~\ref{sec:rem-ledger}, and on Convention~\ref{conv:avg}, the reading of Ba{\l}aban's averaging operator that Theorem~\ref{thm:L046} uses.

\begin{mainthm}[compactness on a finite torus]\label{thm:L003}
Let $G=SU(2)$, $d=4$, $L$ odd and large, and $\ell\in2\mathbb N$ fixed. Let $g_0=g_0(\varepsilon,g)$ be a shooting solution whose effective couplings satisfy $g_k\in\,]0,\gamma]$ for $0\le k\le K$ and $g_K=g$. Assume the uniform bound (H-U) (H-116) of Definition~\ref{def:H116}. Then the lattice Schwinger functions of the centred plaquette density are relatively compact as $\varepsilon=L^{-K}\to0$, and every limit family satisfies the torus forms of (OS0)--(OS3).
\end{mainthm}
This is Theorem~\ref{thm:L013}, proved in Section~\ref{sec:torus}. Theorem~\ref{thm:L003} does not assert that a limit is nonzero. Nontriviality is the separate Theorem~\ref{thm:L014}, which assumes the hypothesis (NT); (NT) is open.

The hypothesis (H-U) is the output of three blocks: the block (B-T2), that is, Theorem~\ref{thm:L002}; the sourced versions of Ba{\l}aban's bounds, block (B-S) (R-035); and the bridge from the source form to (H-U), block (B-N1) (R-045). If (H-U) is taken from these blocks and not assumed, the chain also uses two further inputs. The first input is the $\kappa\to\kappa/2$ decay recovery of the sourced small-field step (R-046). We reduce it, by $L/2$-rescaling, to the large-field T-step rate loss (R-020), which Ba{\l}aban does not print. The second input is the volume hypothesis (Vol), $\ell>22M$ (Section~\ref{sec:lattice}). The dependency table (Section~\ref{sec:deptable}) lists the hypotheses of each block. Section~\ref{sec:remains} lists what remains open.

The sourced step also uses the unit-scale condition (U) (K-024), $\rho_1\le c_*r_0/M^a$ at $\lambda_K=1$. This condition is a choice of the source-gradient constant, made once after $r_0$ and before $K$. It is not a hypothesis on test functions. The bridge (B-N1) meets it by rescaling, at a cost $C^n$. This condition (U) is different from the Route B hypothesis (U).

\emph{Inputs from Ba{\l}aban's papers.} Section~\ref{sec:rem-ledger} lists the statements of Ba{\l}aban and the external results that the two theorems use. Theorem~\ref{thm:L002} uses 4 Ba{\l}aban-level statements and one external result. Theorem~\ref{thm:L003} uses 21 Ba{\l}aban-level statements and four external results. Part (i) of Theorem~\ref{thm:L002}, that is, Theorem 2w, uses one of these statements: R-006, Ba{\l}aban's inductive assumption (H0), which includes his assertion that the whole-lattice limit of \cite[(1.22)]{balaban109} exists. We reduce the use of R-006 in Theorem 2w to named inputs (Section~\ref{sec:src-r5}). The external input of Theorem 2w is the constant-field log coefficient of Savvidy \cite{Savvidy2019}, used in (F). The other inputs of Theorem 2w are Ba{\l}aban's printed material, results proved here relative to listed inputs, conventions, readings, scope restrictions, and conditions on constants. The numerics of Section~\ref{sec:numr8} leave one open question about the single-step $\beta_1$; it does not enter (F).

Of the 23 Ba{\l}aban-level statements of Section~\ref{sec:rem-ledger}, Ba{\l}aban states 16 without proof: 3 of the 4 behind Theorem~\ref{thm:L002}, and 14 of the 21 behind Theorem~\ref{thm:L003}. One more statement, R-021, is a set of printed statements of CMP 99, partly with proof sketches; it contains the printed inputs of the condition (C) (H-057). Each other statement that Ba{\l}aban states without proof, and each statement that he uses without printing it, has a proof here relative to listed inputs or a reduction to named inputs (Sections~\ref{sec:src-r5}, \ref{sec:src-r6}, \ref{sec:src-r7}, and~\ref{sec:rem-ledger}). This includes the rows R-022 and R-023 (CMP 109 pp.~253, 275), whose reductions leave open lemmas, and the row R-024 (the multiscale geometry of CMP 119 pp.~275--276), whose sourced extension we prove relative to listed inputs. The condition (C) itself is reduced to named inputs (Proposition~\ref{prop:L038}), and the volume hypothesis (Vol) discharges its torus-side hypothesis. The lemma (L-LL) (H-137) is Theorem~\ref{thm:L040}, and the bound (S$_\sigma$) (H-168) is Proposition~\ref{prop:L042}.

The propagator and covariance bounds R-001--R-005 reduce to published theorems of CMP 99 and CMP 102. Their joint residue has five parts:
\begin{itemize}
\item Ba{\l}aban's walk expansions of the final propagators, through R-002 (CMP 99 Theorems 3.7 and 3.9 are written out as Theorems~\ref{thm:L045} and~\ref{thm:L044});
\item the complex walk bound (W$^c$) in the exponential gauge, and the decay of $C^*\Delta_kC$ at real $U$ (R-008) under R-003;
\item the regular-class identification R-018;
\item the printed inputs R-021--R-024;
\item the conventions K-044 and K-046.
\end{itemize}

What the paper does not give: literal (0.31) at every $K$ (Section~\ref{sec:literal}), the limit of $\beta_j(0)$ and a rate (Section~\ref{sec:convergence}), uniqueness of the limit, infinite volume, rotation invariance, nontriviality, and a mass gap (Section~\ref{sec:notgiven}). For infinite volume the paper records one conditional statement only, Theorem~\ref{thm:L196}, which is outside both targets. It gives compactness and OS0--OS3 for a joint limit $(\varepsilon,\ell)\to(0,\infty)$, under the hypotheses of Theorem~\ref{thm:L003} at every $\ell$ and two named $\ell$-uniformity lemmas. Theorem~\ref{thm:L003} is a compactness result. Under footnote~2 of Jaffe and Witten \cite{jaffewitten} it is not an admissible solution of the Millennium problem (Section~\ref{sec:exists}).

\subsection{Conventions of this paper}\label{sec:status}

\begin{itemize}
\item \emph{Citations of Ba{\l}aban.} We cite Ba{\l}aban's papers by journal volume and page, for example CMP 109 p.~259. When Ba{\l}aban states a result without proof, we say so and give the page.
\item \emph{Readings.} Where a printed passage admits more than one reading, we fix one reading as a numbered Convention and state the competing reading.
\item \emph{Relative proofs.} ``Proved relative to listed inputs'' means that the proof is complete given the named inputs, and that these inputs are hypotheses of the statement. A ``reduction'' of a statement $X$ to inputs $Y_1,\dots,Y_n$ is a proof of $Y_1\wedge\dots\wedge Y_n\Rightarrow X$; the $Y_i$ stay open unless a reference is given.
\item \emph{Open items.} A statement for which we have no proof appears as a hypothesis, an Open Problem, or a Conjecture. A plausibility argument is labelled as such.
\item \emph{Computations.} A computer-assisted proof uses interval or ball arithmetic with outward rounding over a continuum of parameters. A point enclosure is a rigorous enclosure at one parameter value. All other numerical values are floating-point values without rounding control; they enter no proof.
\item \emph{Lean.} ``Formalized in Lean'' means formalized in Lean 4 with Mathlib, with no \texttt{sorry}, and with the axioms propext, Classical.choice and Quot.sound only (Appendix~\ref{app:lean}).
\item \emph{Project notes.} Some results cite a note of the repository as \file{NAME.md}. A section sign after a note name (for example \file{KAPPA_BOUNDED} \S4) refers to a section of that note, not of this paper. The appendix on provenance lists, for each result, the note that contains the full derivation.
\end{itemize}

\subsection{Relation to other work}\label{sec:independence}

Two recent works overlap with this paper.

Related work 1. E. Sussman and L. Alp\"oge (with Claude), ``The Yang--Mills continuum limit along Ba{\l}aban's renormalization group: a complete account of Theorem 0'' \cite{SA26-T0}, and ``SU(2) Yang--Mills theory in four dimensions via Ba{\l}aban's renormalisation group: construction, axioms and non-triviality'' (draft dated 6 October 2026) \cite{SA26-YM4}. Repository: \url{https://github.com/ethanwsussman/Yang--Mills-LLM-experiment}, commit \texttt{1b92bd0347e3b8763cd8ecc1f96dc3c2a92dd310} (2026-10-27), accessed 2026-10-27. The README states: ``This repository contains work-in-progress, written largely by LLM. Its content has not been satisfactorily vetted. All claims should be understood as tentative, pending verification.''

Related work 2. OpenAI, ``The canonical massive continuum limit of the two-dimensional O(3) model'' (October 4, 2026), \texttt{github.com/openai/math}, family 215 \cite{OAI26-O3}.

The overlaps are as follows.
\begin{enumerate}
\item Sussman and Alp\"oge prove a weak two-sided flow with an additive constant. They do not need per-step positivity, they use a first-passage cutoff, and they use only the containment $g_k\le\gamma$. This is the structure of Theorem 2w.
\item The observation that the weak form of Theorem 2 needs only bounded $\kappa_k$, not their convergence, is theirs (the structure of their Lemma 4.19). It entered this work through a summary of their repository, and Lemma DT and Theorem K use it.
\item The result form of Theorem K (L-009), $|S_k-ck\log L|\le\kappa_*$, matches their Lemma 4.2(b), $|\sum_l\beta^0_l-kc_0|\le c_2$. Our proof of Theorem K (Section~\ref{sec:routeK}) is derived from Ba{\l}aban's papers.
\item Lemma DS (L-004), the shooting argument without early positivity, and the cross-scale sweep of Section~\ref{sec:sweep} were derived without use of their work.
\item Route B's doubling comparison follows the template of \cite[Lemma 6.1]{OAI26-O3}. The idea of a coupling band in place of per-step positivity also appears there.
\end{enumerate}
Their public, dated accounts predate this one.

\subsection{Organization}

Section~\ref{sec:setup} fixes the notation for both parts of the paper. Section~\ref{sec:flow} proves the reduction of Theorem 2 (Theorem~\ref{thm:L002}). Section~\ref{sec:torus} proves the compactness theorem under the hypotheses of the named blocks (Theorem~\ref{thm:L003}) and contains the dependency table. Section~\ref{sec:sources} treats the inputs of the sources block. Section~\ref{sec:numerics} collects the numerical results and the computer-assisted proof of (H2). Section~\ref{sec:remains} lists what remains. Infinite volume, uniqueness, clustering, and the mass gap are outside the scope of this paper. Appendix~\ref{app:lean} lists the Lean formalization, Appendix~\ref{app:ledger} indexes the results, and Appendix~\ref{app:labels} maps the numbering of this paper to the names used in the project notes.
